% Standalone problem document, generated from the adjacent README.md.
% Compile with XeLaTeX. Mathematical content is maintained in Markdown.
\documentclass[11pt,a4paper]{article}
\usepackage[fontsize=10.5pt]{scrextend}
\usepackage[margin=22mm,headheight=15pt,headsep=9mm,footskip=11mm]{geometry}
\usepackage{fontspec}
\setmainfont{texgyrepagella}[Extension=.otf,UprightFont=*-regular,BoldFont=*-bold,ItalicFont=*-italic,BoldItalicFont=*-bolditalic]
\setsansfont{texgyreheros}[Extension=.otf,UprightFont=*-regular,BoldFont=*-bold,ItalicFont=*-italic,BoldItalicFont=*-bolditalic]
\setmonofont{texgyrecursor}[Extension=.otf,UprightFont=*-regular,BoldFont=*-bold,ItalicFont=*-italic,BoldItalicFont=*-bolditalic,Scale=MatchLowercase]
\usepackage{amsmath,amssymb}
\usepackage{unicode-math}
\setmathfont{texgyrepagella-math.otf}
\usepackage{xcolor,microtype,enumitem,needspace,fancyhdr}
\usepackage{bookmark}
\definecolor{ink}{HTML}{17344B}
\definecolor{accent}{HTML}{197D80}
\definecolor{muted}{HTML}{596773}
\hypersetup{colorlinks=true,linkcolor=accent,urlcolor=accent,pdftitle={TR-26 - Degrees
of the two Rayleigh--Ritz discriminant parts for rational normal
curves},pdfauthor={Open Problems in Numerical Linear Algebra}}
\urlstyle{same}
\setlength{\parindent}{0pt}
\setlength{\parskip}{5pt plus 1pt minus 1pt}
\linespread{1.04}
\setlength{\emergencystretch}{3em}
\widowpenalty=10000
\clubpenalty=10000
\displaywidowpenalty=10000
\allowdisplaybreaks[1]
\setlist{nosep,leftmargin=1.5em}
\providecommand{\tightlist}{\setlength{\itemsep}{0pt}\setlength{\parskip}{0pt}}
\providecommand{\pandocbounded}[1]{#1}
\pagestyle{fancy}
\fancyhf{}
\fancyhead[L]{\sffamily\footnotesize\color{muted}OPEN PROBLEMS IN NUMERICAL LINEAR ALGEBRA}
\fancyhead[R]{\sffamily\bfseries\footnotesize\color{accent}TR-26}
\fancyfoot[L]{\parbox[b]{0.9\textwidth}{\sffamily\fontsize{8}{10}\selectfont\color{muted}Editorial ratings; a bounded search cannot certify that a problem remains unsolved.\\Literature check: 2026-09-10}}
\fancyfoot[R]{\sffamily\footnotesize\color{muted}\thepage}
\renewcommand{\headrulewidth}{0.3pt}
\renewcommand{\headrule}{\hbox to\headwidth{\color{accent}\leaders\hrule height \headrulewidth\hfill}}
\makeatletter
\renewcommand{\section}{\Needspace{5\baselineskip}\@startsection{section}{1}{0pt}{12pt plus 2pt minus 2pt}{4pt}{\sffamily\bfseries\large\color{ink}}}
\renewcommand{\subsection}{\Needspace{4\baselineskip}\@startsection{subsection}{2}{0pt}{9pt plus 1pt minus 1pt}{3pt}{\sffamily\bfseries\normalsize\color{ink}}}
\makeatother
\setcounter{secnumdepth}{0}
\begin{document}
{\sffamily\small\color{muted}Tensor computations\par}
\vspace{3pt}
{\raggedright\hyphenpenalty=10000\exhyphenpenalty=10000\sffamily\bfseries\fontsize{21}{24}\selectfont\color{ink}Degrees
of the two Rayleigh--Ritz discriminant parts for rational normal
curves\par}
\vspace{7pt}
{\sffamily\small\textbf{Difficulty:} challenging\quad\textbf{Importance:} interesting
to specialist\par}
{\sffamily\small\bfseries\color{accent}Status: Partially resolved\par}
\vspace{4pt}
{\color{accent}\hrule height 0.8pt}
\vspace{6pt}
\textbf{Rating rationale:} Separating the isotropic and nonisotropic
degree contributions requires substantial algebraic analysis beyond the
known total. The result mainly interests specialists in conditioning of
constrained eigenvalue problems.

\subsection{Problem statement}\label{problem-statement}

For every integer \(d\geq2\), use the specific rational normal curve \[
C_d=\{[a^d:a^{d-1}b:\cdots:b^d]:[a:b]\in\mathbb P^1\}
\subset\mathbb P^d.
\] For a complex symmetric matrix \(H\) of order \(d+1\), consider
\(\rho_H(z)=z^\top Hz/(z^\top z)\) on the part of \(C_d\) where
\(z^\top z\ne0\). Transposes here do not conjugate.

Let \(\mathcal R\) be the Zariski closure in
\(C_d\times\mathbb P(\operatorname{Sym}_{d+1}(\mathbb C))\) of all pairs
\(([z],[H])\) for which \([z]\) is a critical point of
\(\rho_H|_{C_d}\). Choose homogeneous generators \(g_i\) of its ideal
and let \[
\mathcal B=\{([z],[H])\in\mathcal R:
\operatorname{rank}(\partial g_i/\partial z_j)\leq d-1\}.
\] Put \(Q=\{([z],[H]):z^\top z=0\}\) and let \(\pi\) project to the
matrix factor. Define \[
\Sigma_{\mathrm{iso}}=\pi(\mathcal B\cap Q),\qquad
\Sigma_{\mathrm{noniso}}=\overline{\pi(\mathcal B\setminus Q)}^{\,\mathrm{Zar}}.
\] Are these projective hypersurfaces with \[
\deg\Sigma_{\mathrm{noniso}}=6(d-1),\qquad
\deg\Sigma_{\mathrm{iso}}=2d?
\] Degree refers to each reduced algebraic set, including all its
components. The question is Conjecture 3.14 in the source. Both parts
belong to one problem; the standard unweighted embedding and the
bilinear quadratic form are fixed.

\subsection{Why it matters}\label{why-it-matters}

These discriminants locate parameter values at which constrained
Rayleigh-quotient critical points collide or cease to form a reduced
finite set. They describe degeneracy in a basic nonlinear analogue of
matrix eigenvalue computations.

\newpage
\subsection{References and status
check}\label{references-and-status-check}

\begin{enumerate}
\def\labelenumi{\arabic{enumi}.}
\tightlist
\item
  V. Borovik, H. Friedman, S. Hoşten, and M. Pfeffer, \emph{Numerical
  Algebraic Geometry for Energy Computations on Tensor Train Varieties},
  \href{https://arxiv.org/pdf/2512.06939v2}{arXiv:2512.06939v2}, §3.2,
  Conjecture 3.14, pp.13--14; the correspondence and ramification
  definitions are on pp.7--13.
\item
  The same paper, Proposition 3.10, proves the total discriminant degree
  \(2(4d-3)\); Examples 3.12--3.13 compute small cases.
\end{enumerate}

\subsubsection{Status check ---
2026-09-10}\label{status-check-2026-09-10}

Checked the June 2026 revision, Proposition 3.10, Examples 3.12--3.13
and Conjecture 3.14, and searched the title and rational-normal-curve
Rayleigh--Ritz discriminants. Example 3.12 gives an explicit symbolic
factorization for \(d=2\), with isotropic degree four and nonisotropic
degree six; this is a proved small-degree case supporting Partially
resolved. Example 3.13 also reports symbolic degree computation for
\(d=3\). The separate numerical checks through \(d=10\) are evidence,
not a general proof. The total-degree theorem alone does not determine
the component split. No later general proof or counterexample was
located.
\end{document}
